\documentclass[10pt,a4paper]{amsart}
\usepackage{amsmath,amssymb,amsthm,graphicx}
\usepackage{verbatim}
\newenvironment{reference}{\comment}{\endcomment}
%\newenvironment{reference}{}{}
\newenvironment{slogan}{\comment}{\endcomment}
\newenvironment{history}{\comment}{\endcomment}

% For commutative diagrams we use Xy-pic
\usepackage[all]{xy}

% We use 2cell for 2-commutative diagrams.
\xyoption{2cell}
% Theorem environments.
%
\theoremstyle{plain}
\newtheorem{theorem}[subsection]{Theorem}
\newtheorem{proposition}[subsection]{Proposition}
\newtheorem{lemma}[subsection]{Lemma}

\theoremstyle{definition}
\newtheorem{definition}[subsection]{Definition}
\newtheorem{example}[subsection]{Example}
\newtheorem{exercise}[subsection]{Exercise}
\newtheorem{situation}[subsection]{Situation}

\theoremstyle{remark}
\newtheorem{remark}[subsection]{Remark}
\newtheorem{remarks}[subsection]{Remarks}

\numberwithin{equation}{subsection}

% Macros
%
\def\lim{\mathop{\mathrm{lim}}\nolimits}
\def\colim{\mathop{\mathrm{colim}}\nolimits}
\def\Spec{\mathop{\mathrm{Spec}}}
\def\Hom{\mathop{\mathrm{Hom}}\nolimits}
\def\Ext{\mathop{\mathrm{Ext}}\nolimits}
\def\SheafHom{\mathop{\mathcal{H}\!\mathit{om}}\nolimits}
\def\SheafExt{\mathop{\mathcal{E}\!\mathit{xt}}\nolimits}
\def\Sch{\mathit{Sch}}
\def\Mor{\mathop{\mathrm{Mor}}\nolimits}
\def\Ob{\mathop{\mathrm{Ob}}\nolimits}
\def\Sh{\mathop{\mathit{Sh}}\nolimits}
\def\NL{\mathop{N\!L}\nolimits}
\def\CH{\mathop{\mathrm{CH}}\nolimits}
\def\proetale{{pro\text{-}\acute{e}tale}}
\def\etale{{\acute{e}tale}}
\def\QCoh{\mathit{QCoh}}
\def\Ker{\mathop{\mathrm{Ker}}}
\def\Im{\mathop{\mathrm{Im}}}
\def\Coker{\mathop{\mathrm{Coker}}}
\def\Coim{\mathop{\mathrm{Coim}}}

% Boxtimes
%
\DeclareMathSymbol{\boxtimes}{\mathbin}{AMSa}{"02}

%
% Macros for moduli stacks/spaces
%
\def\QCohstack{\mathcal{QC}\!\mathit{oh}}
\def\Cohstack{\mathcal{C}\!\mathit{oh}}
\def\Spacesstack{\mathcal{S}\!\mathit{paces}}
\def\Quotfunctor{\mathrm{Quot}}
\def\Hilbfunctor{\mathrm{Hilb}}
\def\Curvesstack{\mathcal{C}\!\mathit{urves}}
\def\Polarizedstack{\mathcal{P}\!\mathit{olarized}}
\def\Complexesstack{\mathcal{C}\!\mathit{omplexes}}
% \Pic is the operator that assigns to X its picard group, usage \Pic(X)
% \Picardstack_{X/B} denotes the Picard stack of X over B
% \Picardfunctor_{X/B} denotes the Picard functor of X over B
\def\Pic{\mathop{\mathrm{Pic}}\nolimits}
\def\Picardstack{\mathcal{P}\!\mathit{ic}}
\def\Picardfunctor{\mathrm{Pic}}
\def\Deformationcategory{\mathcal{D}\!\mathit{ef}}
\title{fppf torsors under a smooth commutative group algebraic space over a field are torsion}
\author{The Stacks Project Authors}
\date{stacks-project:examples.tex:local-snapshot-20261012; GNU FDL 1.2 or later}
\begin{document}
\maketitle
\noindent\textit{Editorial note.} All delivered mathematical text is the author's own, copied from the pinned Stacks Project chapter file examples.tex; no statement or proof step was rewritten, reordered or removed.\par\medskip
\noindent\textit{Editorial note.} The counted proof is the author's complete argument: The author leaves one verification to the reader ('the reader can verify that the scheme on the left has a K-rational point'); recorded as author's own style, content verbatim, no rewrite.\par\medskip
\noindent\textit{Editorial note.} The label lemma-torsors-over-field-torsion names the counted statement, and the proof environment immediately following it in the source is delivered verbatim as its proof.\par\medskip
\noindent\textit{Editorial note.} The counted statement is carried by the author's own lemma environment in the chapter source.\par\medskip
\noindent\textit{Allowed standard prerequisites (editorial).}\begin{itemize}
\item fppf cohomology H\textasciicircum{}1 and G-torsors over Spec(K) for G a smooth commutative group algebraic space over a field K
\item the torsor [n](P) whose class is n times the class of P in H\textasciicircum{}1\_fppf(Spec(K), G)
\item smooth schemes/spaces have points over finite separable extensions; quotients by finite constant groups
\end{itemize}\medskip
\begin{lemma}
\label{lemma-torsors-over-field-torsion}
Let $G$ be a smooth commutative group algebraic space over a field $K$.
Then $H^1_{fppf}(\Spec(K), G)$ is torsion.
\end{lemma}

\begin{proof}
Every $G$-torsor $P$ over $\Spec(K)$ is smooth over $K$ as a form of $G$.
Hence $P$ has a point over a finite separable extension $L/K$.
Say $[L : K] = n$. Let $[n](P)$ denote the $G$-torsor whose class is $n$
times the class of $P$ in $H^1_{fppf}(\Spec(K), G)$. There is a canonical
morphism
$$
P \times_{\Spec(K)} \ldots \times_{\Spec(K)} P \to [n](P)
$$
of algebraic spaces over $K$. This morphism is symmetric as
$G$ is abelian. Hence it factors through the quotient
$$
(P \times_{\Spec(K)} \ldots \times_{\Spec(K)} P)/S_n
$$
On the other hand, the morphism $\Spec(L) \to P$ defines a morphism
$$
(\Spec(L) \times_{\Spec(K)} \ldots \times_{\Spec(K)} \Spec(L))/S_n
\longrightarrow (P \times_{\Spec(K)} \ldots \times_{\Spec(K)} P)/S_n
$$
and the reader can verify that the scheme on the left has a $K$-rational
point. Thus we see that $[n](P)$ is the trivial torsor.
\end{proof}

\end{document}
